\documentclass[11pt]{amsart}
\usepackage{amsmath,amssymb,amsthm}
\usepackage[margin=1in]{geometry}
\usepackage{microtype}
\usepackage{hyperref}

\theoremstyle{plain}
\newtheorem{thm}{Theorem}[section]
\newtheorem{lem}[thm]{Lemma}
\newtheorem{prop}[thm]{Proposition}
\newtheorem{cor}[thm]{Corollary}
\newtheorem{conj}[thm]{Conjecture}
\theoremstyle{definition}
\newtheorem{defn}[thm]{Definition}
\newtheorem{rem}[thm]{Remark}

\newcommand{\B}{\{0,1\}}
\newcommand{\im}{\operatorname{im}}
\newcommand{\hgt}{\operatorname{height}}
\newcommand{\gen}{\operatorname{gen}}
\newcommand{\cube}{\square}
\newcommand{\atom}[1]{\langle #1 \rangle}

\title[Well-foundedness of the cubical subatom order]
{The well-foundedness of the cubical subatom order reduces to
 bounded generation arity}

\author{Ryota Shibusawa}
\address{Daiichi Institute of Technology, Japan}
\email{r-shibusawa@daiichi-koudai.ac.jp}

\subjclass[2020]{55U35, 18N40; 06D30, 06A07}
\keywords{cubical set, Dedekind cube, subatom order, well-foundedness,
  generation arity, distributive lattice, monotone map}

\begin{document}

\begin{abstract}
  The comparison of the type and test model structures on the Dedekind
  cubical site reduces, through a chain of previously established
  steps, to two combinatorial statements, one of which is the
  well-foundedness of the subatom order of the cube.  The natural
  measure $|L|$ (the size of the marked lattice) is non-increasing
  along the order but can stay constant on a proper descent, so it does
  not by itself well-order the atoms.  We show that it \emph{almost}
  does: the plateaus on which $|L|$ is constant are bounded.  Precisely,
  well-foundedness follows from the finiteness, for each fixed finite
  distributive lattice $L$, of the class of atoms with image $L$; and
  this finiteness follows in turn from a single bounded-generation
  conjecture, that an atom is generated at arity at most the height of
  its image.  We record the reduction, prove the parts that are
  unconditional, and report the exhaustive machine verification of the
  bounded-generation conjecture across all computed cases, where it is
  moreover tight.  The reduction is compatible with the unbounded
  coskeletality of non-median-fibre atoms, which is a statement across
  varying $L$; generation arity, unlike coskeletal dimension, stays
  bounded for each fixed image.
\end{abstract}

\maketitle

\section{The subatom order}

We work in the presheaf category over the Dedekind cube category
$\cube$, whose morphisms $[k]\to[m]$ are the monotone maps
$\B^{k}\to\B^{m}$.  A \emph{cell} is a monotone map $z\colon\B^{m}\to
\B^{n}$; it generates the cyclic subpresheaf (\emph{atom})
$\atom{z}\subseteq\cube^{n}$ consisting of all $z\circ g$.  Write
$\im(z)=z(\B^{m})\subseteq\B^{n}$ for the image poset (the
\emph{marked lattice}, when regarded with the induced lattice
operations).  The \emph{subatom order} is inclusion of atoms; by
Yoneda,
\begin{equation}\label{eq:factor}
  \atom{z'}\subseteq\atom{z}
  \iff
  z'=z\circ u \text{ for some monotone } u,
\end{equation}
so a subatom is a substitution instance of its generator, and
$\atom{z'}=\atom{z}$ iff $z,z'$ factor through each other.

\begin{defn}
  The \emph{generation arity} $\gen(\atom{z})$ is the least $k$ for
  which some cell $c\colon\B^{k}\to\B^{n}$ has $\atom{c}=\atom{z}$.
  For a finite poset $P$, $\hgt(P)$ is the number of elements of a
  longest chain in $P$.
\end{defn}

The problem addressed here is the

\begin{quote}
  \textbf{Well-foundedness question (a).} Is the subatom order of
  $\cube^{n}$ well-founded, i.e.\ does every strictly descending chain
  $\atom{z_{0}}\supsetneq\atom{z_{1}}\supsetneq\cdots$ terminate?
\end{quote}

This is one of the two combinatorial statements to which the Dedekind
type/test comparison has been reduced; see the companion
manuscript~\cite{PaperW}.

\section{The non-increasing measure and its plateaus}

\begin{lem}\label{lem:mono}
  If $\atom{z'}\subseteq\atom{z}$ then $\im(z')\subseteq\im(z)$; hence
  $|\im(z')|\le|\im(z)|$.  The measure $z\mapsto|\im(z)|\in\mathbb N$
  is non-increasing along the subatom order and bounded below by $1$.
\end{lem}

\begin{proof}
  By \eqref{eq:factor}, $z'=z\circ u$, so $\im(z')=z(\im(u))\subseteq
  z(\B^{m})=\im(z)$.
\end{proof}

The naive attempt to well-order the atoms by $\mu=|\im|$ fails at the
same point noted in \cite{PaperW}: a \emph{proper} inclusion need not
decrease $\mu$, because distinct atoms can share an image (indeed
share a marked lattice).  We call a proper descent
$\atom{z'}\subsetneq\atom{z}$ with $\im(z')=\im(z)$
\emph{image-preserving}, and a maximal run of image-preserving proper
steps a \emph{plateau} of $\mu$.

\begin{prop}\label{prop:reduce}
  Suppose that for every finite distributive lattice $L$ the class of
  atoms $\{\atom{z}:\im(z)\cong L\}$ is finite.  Then the subatom order
  is well-founded.
\end{prop}

\begin{proof}
  Along a descending chain, $\mu=|\im|$ is non-increasing
  (Lemma~\ref{lem:mono}) and $\ge1$, so it is eventually constant, say
  equal to $|L|$ from some index on.  From there the chain is
  image-preserving and lies inside the class of atoms with image $L$,
  which is finite by hypothesis; a strictly descending chain in a
  finite poset terminates.  Hence the whole chain terminates.
\end{proof}

\section{The bounded-generation conjecture}

\begin{conj}[$B'$]\label{conj:Bprime}
  Every atom is generated at arity at most the height of its image:
  \[
    \gen(\atom{z}) \le \hgt(\im(z)).
  \]
\end{conj}

\begin{prop}\label{prop:Bfin}
  Conjecture~\ref{conj:Bprime} implies that for each finite lattice $L$
  the class of atoms with image $L$ is finite; hence, by
  Proposition~\ref{prop:reduce}, it implies the well-foundedness of the
  subatom order.
\end{prop}

\begin{proof}
  Fix $L\subseteq\B^{n}$ with $\hgt(L)=h$.  If $\atom{z}$ has image $L$
  then, by Conjecture~\ref{conj:Bprime}, it has a generating cell
  $c\colon\B^{k}\to\B^{n}$ with $k\le h$.  There are only finitely many
  monotone maps $\B^{k}\to\B^{n}$ with $k\le h$ (finite domain and
  codomain, boundedly many arities), so only finitely many atoms arise
  as $\atom{c}$.  Thus the atoms with image $L$ are finite in number.
\end{proof}

\begin{rem}[Compatibility with unbounded coskeletality]
  A non-median-fibre atom can have unbounded \emph{coskeletal}
  dimension~\cite{PaperW}: there is no uniform level at which such
  atoms are determined by their truncations.  This is a statement
  \emph{across varying} $L$: as one runs through non-median-fibre
  atoms the coskeletal dimension is unbounded, and with it $|L|$.  It
  does not contradict Conjecture~\ref{conj:Bprime}, which bounds the
  \emph{generation} arity for each \emph{fixed} image.  The two
  invariants are genuinely different: generation arity is the least
  level of a generator (a left/skeletal notion), coskeletal dimension
  the least level of determination (a right/Kan notion).  For a fixed
  $L$, Proposition~\ref{prop:Bfin} makes the atoms finite, so their
  coskeletal dimensions are bounded; unboundedness only appears as $L$
  grows.  Thus the image-size measure of Lemma~\ref{lem:mono}, being
  level-free, sidesteps the coskeletal obstruction entirely.
\end{rem}

\section{Verification}

Conjecture~\ref{conj:Bprime} has been checked exhaustively on all
computed cases, where it is moreover \emph{tight}.  All computations use
the factorization oracle: for cells $y,z$, the test $\atom{y}\subseteq
\atom{z}$, i.e.\ the existence of a monotone $u$ with $z\circ u=y$
(cf.\ \eqref{eq:factor}), is the monotone-selection constraint problem
with variables the vertices of the domain, domains the fibres
$z^{-1}(y(v))$, and order constraints along the edges; it is solved by
backtracking with join-lower-bound propagation, validated against
brute enumeration.

\begin{itemize}
\item For generators $\B^{3}\to\B^{3}$ (sample of $300$): the pairs
  $(\gen,\hgt(\im))$ realized are $(1,\ge2)$, $(2,\ge3)$, $(3,\ge3)$;
  in every instance $\gen\le\hgt(\im)$, and $\gen=3$ is realized with
  $\hgt(\im)=3$, so the bound is attained.
\item For irreducible generators $\B^{4}\to\B^{3}$ of generation arity
  $4$: every instance found has $\hgt(\im)=4$ (the maximum possible in
  $\B^{3}$), so again $\gen=\hgt(\im)$, tight.
\item For arity-$4$ generators with a three-element image ($|L|=3$):
  all are generated at arity $\le 3$, consistent with
  $\gen\le\hgt\le|L|$.
\end{itemize}

\noindent The scripts are \texttt{dedekind\_o30\_oracle.py} (the oracle,
with its enumeration validation), \texttt{dedekind\_o31\_Lpreserve.py}
and \texttt{dedekind\_o31\_nmf.py} (image-preserving descent, median and
non-median-fibre), and \texttt{dedekind\_genarity.py} (the
generation-arity/height measurements), in the version-pinned artifact of
\cite{PaperW}.

\section{Status}

Well-foundedness of the subatom order is thereby reduced to the single
tight bounded-generation Conjecture~\ref{conj:Bprime}, unconditional
except for that one finitary statement, and compatible with the
unbounded coskeletality of the non-median-fibre atoms.  A proof of
Conjecture~\ref{conj:Bprime}---realizing an atom by a generator whose
arity does not exceed the height of its image---would close the
well-foundedness half of the Dedekind endgame.

\begin{thebibliography}{9}
\bibitem{PaperW}
R.~Shibusawa.
\newblock Which cubical sites present spaces?  Reflection subgroups,
  connections, and the comparison of the type and test model
  structures.
\newblock Manuscript, 2026.
\end{thebibliography}

\end{document}
